% GENERATED FILE. Edit canonical lessons, then run npm run generate:books.
% canonical-source-hash: fnv1a64:eaed80977b5d9f07
% canonical-lessons: ES-C775-nevera, ES-C775-cafetera, ES-C775-lata, ES-C775-trozo, ES-C775-olor, ES-R775-repaso-768-775

\chapter{En mi cocina}
\label{ch:es-en-mi-cocina}
\hlchaptermodality{775}

\begin{chapteropening}
\textbf{By the end of this chapter:} \emph{I can describe my kitchen and what I find there: the fridge, the coffee maker, a tin, a piece of something, and the smell.}

Complete the last lesson of chapter 775: I can describe my kitchen and what I find there: the fridge, the coffee maker, a tin, a piece of something, and the smell.
\end{chapteropening}

\section[la nevera]{la nevera — the fridge}

\label{lesson:ES-C775-nevera}



\begin{quote}
Before the new one: say the Spanish for the mirror, then the Spanish for the blind.
\end{quote}

\subsection*{You'll want to know: la nevera}
\textbf{la nevera} — \textquotedblleft{}the fridge\textquotedblright{}.

\begin{quote}
La nevera está casi vacía.
\end{quote}

\begin{quote}
\emph{The fridge is almost empty.}
\end{quote}

\emph{El frigorífico} is also common in Spain; in parts of Latin America people say \emph{el refrigerador} or \emph{la heladera}.

\begin{cousinweb}
\emph{nevera} comes from Latin \emph{nivārius} 'of snow' (EtymDB; Lewis \& Short), built on \emph{nix}, \emph{nivis} 'snow', which also gave \emph{nieve}. It was first a place where snow was stored to keep things cold.
\end{cousinweb}

\begin{grammarlens}[title={What you've built}]
One more word for telling, explaining and describing at B1.
\end{grammarlens}

\subsection*{Your turn}
\begin{itemize}
\raggedright
  \item \emph{Say it:} \emph{la nevera}
  \item \emph{Say it:} \emph{la nevera}, once more
  \item \emph{Say it:} a sentence of your own with \emph{la nevera}
  \item \emph{Recall:} say the Spanish for the mirror, then the Spanish for the blind, then say \emph{la nevera} again
\end{itemize}

\begin{tcolorbox}[breakable,colback=teal!4,colframe=teal!35!black,title={Before you move on}]
What does la nevera mean? (\textquotedblleft{}The fridge\textquotedblright{}.) Say it once more.
\end{tcolorbox}
\section[la cafetera]{la cafetera — the coffee maker, the coffee pot}

\label{lesson:ES-C775-cafetera}



\begin{quote}
Before the new one: say the Spanish for the blind, then the Spanish for the fridge.
\end{quote}

\subsection*{You'll want to know: la cafetera}
\textbf{la cafetera} — \textquotedblleft{}the coffee maker, the coffee pot\textquotedblright{}.

\begin{quote}
Por la mañana la cafetera hace mucho ruido.
\end{quote}

\begin{quote}
\emph{In the morning the coffee maker makes a lot of noise.}
\end{quote}

\begin{cousinweb}
\emph{cafetera} comes from French \emph{cafetière} 'coffee pot' (EtymDB), built on \emph{café}, which reached Europe through Ottoman Turkish \emph{kahve} from Arabic \emph{qahwa}.
\end{cousinweb}

\begin{grammarlens}[title={What you've built}]
One more word for telling, explaining and describing at B1.
\end{grammarlens}

\subsection*{Your turn}
\begin{itemize}
\raggedright
  \item \emph{Say it:} \emph{la cafetera}
  \item \emph{Say it:} \emph{la cafetera}, once more
  \item \emph{Say it:} a sentence of your own with \emph{la cafetera}
  \item \emph{Recall:} say the Spanish for the blind, then the Spanish for the fridge, then say \emph{la cafetera} again
\end{itemize}

\begin{tcolorbox}[breakable,colback=teal!4,colframe=teal!35!black,title={Before you move on}]
What does la cafetera mean? (\textquotedblleft{}The coffee maker, the coffee pot\textquotedblright{}.) Say it once more.
\end{tcolorbox}
\section[la lata]{la lata — the tin, the can; a nuisance}

\label{lesson:ES-C775-lata}



\begin{quote}
Before the new one: say the Spanish for the fridge, then the Spanish for the coffee maker.
\end{quote}

\subsection*{You'll want to know: la lata}
\textbf{la lata} — \textquotedblleft{}the tin, the can; a nuisance\textquotedblright{}.

\begin{quote}
En el armario hay latas de atún.
\end{quote}

\begin{quote}
\emph{In the cupboard there are tins of tuna.}
\end{quote}

\emph{¡Qué lata!} — what a pain! \emph{Dar la lata} — to pester.

\begin{cousinweb}
\emph{lata} is usually traced to Late Latin \emph{latta} 'a lath, a thin board', probably of Germanic origin. EtymDB traces its French cousin \emph{latte} to Frankish \emph{*latta}, the same Germanic word as English \emph{lath}. The sense moved from a thin strip of wood to a thin sheet of tin.
\end{cousinweb}

\begin{grammarlens}[title={What you've built}]
One more word for telling, explaining and describing at B1.
\end{grammarlens}

\subsection*{Your turn}
\begin{itemize}
\raggedright
  \item \emph{Say it:} \emph{la lata}
  \item \emph{Say it:} \emph{la lata}, once more
  \item \emph{Say it:} a sentence of your own with \emph{la lata}
  \item \emph{Recall:} say the Spanish for the fridge, then the Spanish for the coffee maker, then say \emph{la lata} again
\end{itemize}

\begin{tcolorbox}[breakable,colback=teal!4,colframe=teal!35!black,title={Before you move on}]
What does la lata mean? (\textquotedblleft{}The tin, the can; a nuisance\textquotedblright{}.) Say it once more.
\end{tcolorbox}
\section[el trozo]{el trozo — the piece, the bit}

\label{lesson:ES-C775-trozo}



\begin{quote}
Before the new one: say the Spanish for the coffee maker, then the Spanish for the tin.
\end{quote}

\subsection*{You'll want to know: el trozo}
\textbf{el trozo} — \textquotedblleft{}the piece, the bit\textquotedblright{}.

\begin{quote}
Me queda un trozo de pastel en la nevera.
\end{quote}

\begin{quote}
\emph{I have a piece of cake left in the fridge.}
\end{quote}

\emph{Un trozo de pan} — a piece of bread.

\begin{cousinweb}
The origin of \emph{trozo} is not settled. One proposal links it to Latin \emph{thyrsus} 'a stalk' (Lewis \& Short), but no account is agreed.
\end{cousinweb}

\begin{grammarlens}[title={What you've built}]
One more word for telling, explaining and describing at B1.
\end{grammarlens}

\subsection*{Your turn}
\begin{itemize}
\raggedright
  \item \emph{Say it:} \emph{el trozo}
  \item \emph{Say it:} \emph{el trozo}, once more
  \item \emph{Say it:} a sentence of your own with \emph{el trozo}
  \item \emph{Recall:} say the Spanish for the coffee maker, then the Spanish for the tin, then say \emph{el trozo} again
\end{itemize}

\begin{tcolorbox}[breakable,colback=teal!4,colframe=teal!35!black,title={Before you move on}]
What does el trozo mean? (\textquotedblleft{}The piece, the bit\textquotedblright{}.) Say it once more.
\end{tcolorbox}
\section[el olor]{el olor — the smell}

\label{lesson:ES-C775-olor}



\begin{quote}
Before the new one: say the Spanish for the tin, then the Spanish for the piece.
\end{quote}

\subsection*{You'll want to know: el olor}
\textbf{el olor} — \textquotedblleft{}the smell\textquotedblright{}.

\begin{quote}
Me encanta el olor del pan recién hecho.
\end{quote}

\begin{quote}
\emph{I love the smell of freshly baked bread.}
\end{quote}

\emph{Olor a} — a smell of: \emph{olor a café}.

\begin{cousinweb}
\emph{olor} is usually traced to Latin \emph{olor} 'a smell' (Lewis \& Short), a rare by-form of \emph{odor} shaped by the verb \emph{olēre} 'to give off a smell'.
\end{cousinweb}

\begin{grammarlens}[title={What you've built}]
One more word for telling, explaining and describing at B1.
\end{grammarlens}

\subsection*{Your turn}
\begin{itemize}
\raggedright
  \item \emph{Say it:} \emph{el olor}
  \item \emph{Say it:} \emph{el olor}, once more
  \item \emph{Say it:} a sentence of your own with \emph{el olor}
  \item \emph{Recall:} say the Spanish for the tin, then the Spanish for the piece, then say \emph{el olor} again
\end{itemize}

\begin{tcolorbox}[breakable,colback=teal!4,colframe=teal!35!black,title={Before you move on}]
What does el olor mean? (\textquotedblleft{}The smell\textquotedblright{}.) Say it once more.
\end{tcolorbox}
\section[(dialogue)]{Casa, cocina y compras}

\label{lesson:ES-R775-repaso-768-775}



\begin{quote}
Say the Spanish for to boil and to fry.
\end{quote}

\subsection*{Your turn}
Say these aloud:

\begin{itemize}
\raggedright
  \item to boil, to fry, to peel, to beat, the ingredient
  \item to drip, the leak, broken down, the plumber, the electrician
  \item the dry cleaner's, the fitting room, the hardware shop, the butcher's, the greengrocer's
  \item the taste, spicy, salty, raw, to smell
  \item the breeze, to blow, the hillside, to slip, the compass
  \item the signal, the adapter, the app, the link, to top up
  \item the basement, the attic, the drawer, the mirror, the blind
  \item the fridge, the coffee maker, the tin, the piece, the smell
\end{itemize}

\begin{tcolorbox}[breakable,colback=teal!4,colframe=teal!35!black,title={Before you move on}]
To boil? (\textbf{hervir}.) To fry? (\textbf{freír}.) To peel? (\textbf{pelar}.) To beat? (\textbf{batir}.) The ingredient? (\textbf{el ingrediente}.) To drip? (\textbf{gotear}.) The leak? (\textbf{la gotera}.) Broken down? (\textbf{averiado}.) The plumber? (\textbf{el fontanero}.) The electrician? (\textbf{el electricista}.) The dry cleaner's? (\textbf{la tintorería}.) The fitting room? (\textbf{el probador}.) The hardware shop? (\textbf{la ferretería}.) The butcher's? (\textbf{la carnicería}.) The greengrocer's? (\textbf{la frutería}.) The taste? (\textbf{el sabor}.) Spicy? (\textbf{picante}.) Salty? (\textbf{salado}.) Raw? (\textbf{crudo}.) To smell? (\textbf{oler}.) The breeze? (\textbf{la brisa}.) To blow? (\textbf{soplar}.) The hillside? (\textbf{la ladera}.) To slip? (\textbf{resbalar}.) The compass? (\textbf{la brújula}.) The signal? (\textbf{la cobertura}.) The adapter? (\textbf{el adaptador}.) The app? (\textbf{la aplicación}.) The link? (\textbf{el enlace}.) To top up? (\textbf{recargar}.) The basement? (\textbf{el sótano}.) The attic? (\textbf{el desván}.) The drawer? (\textbf{el cajón}.) The mirror? (\textbf{el espejo}.) The blind? (\textbf{la persiana}.) The fridge? (\textbf{la nevera}.) The coffee maker? (\textbf{la cafetera}.) The tin? (\textbf{la lata}.) The piece? (\textbf{el trozo}.) The smell? (\textbf{el olor}.)
\end{tcolorbox}
